\section{The inductive framework}\label{proj:sec:induction}

All varieties in this section are projective over \(\C\).
We first isolate the established minimal-model results that will be
used, and explain exactly where the new signed-representative theorem
enters the induction.

For an effective real boundary \(\Delta\), a good log minimal model of
\((X,\Delta)\) is a log minimal model on which the adjoint divisor is
semiample. For real divisors, semiampleness means real linear equivalence
to the pullback of an ample real divisor by a contraction. We allow the
usual definition of a log minimal model in which extracted prime
divisors are included in its boundary with coefficient one.
For a pseudo-effective divisor \(A\), we use
\(\kappa_\sigma(X,A)\) for Nakayama's numerical dimension. When \(A\)
is nef this agrees with the intersection-theoretic numerical dimension
\(\nu(X,A)\).

Fix an integer \(n\geq 1\).
\begin{assumption}[Lower-dimensional good models]\label{proj:mmp:lower-models}
Every projective log canonical pair of dimension less than \(n\),
with real boundary and pseudo-effective adjoint divisor, has a good
log minimal model.
\end{assumption}

The dimension-zero case is immediate. Our inductive step will first
prove smooth nonvanishing in dimension \(n\) under
Assumption~\ref{proj:mmp:lower-models}, and then apply
Proposition~\ref{proj:prop:inductive-reduction} below.
The propositions on boundary restrictions and special termination
preceding that reduction do not assume smooth nonvanishing in
dimension \(n\).

\subsection{Comparison and boundary restrictions}

\begin{lemma}[Comparison with a nef model]\label{proj:mmp:comparison}
Let \((X,\Delta)\) be log canonical and let \((X',\Delta')\)
be a log minimal model. On a common resolution
\[
 X \xleftarrow{u} W \xrightarrow{v} X'
\]
there is an effective \(v\)-exceptional divisor \(F\) such that
\begin{equation}\label{proj:mmp:comparison-equation}
 u^*(K_X+\Delta)=v^*(K_{X'}+\Delta')+F.
\end{equation}
If \(K_X+\Delta\) is nef, then \(F=0\). Consequently a nef
rational adjoint is semiample whenever it has a good log minimal
model. If \((X,\Delta)\) is klt, its log minimal model is klt and
extracts no divisors.
\end{lemma}

\begin{proof}
Put \(F=u^*(K_X+\Delta)-v^*(K_{X'}+\Delta')\).
Discrepancy improvement on divisors of \(X\) gives \(u_*F\geq 0\),
while \(-F\) is \(u\)-nef because \(K_{X'}+\Delta'\) is nef.
The negativity lemma therefore gives \(F\geq 0\).
At a prime of \(W\) that is not \(v\)-exceptional, the coefficient
of \(F\) is zero unless that prime is extracted on \(X'\).
In the latter case it equals the negative of its log discrepancy
over \((X,\Delta)\). It is therefore nonpositive, and hence zero.
Thus \(F\) is \(v\)-exceptional. If the source adjoint is nef,
then \(F\) is also \(v\)-nef, so the negativity lemma gives \(F=0\).

Equality of the pullbacks transfers semiampleness, since a proper
birational morphism onto a normal variety has direct image of its
structure sheaf equal to the structure sheaf. More explicitly,
global sections of a Cartier multiple and its pullback coincide;
generation upstairs implies generation downstairs.
Finally, for a klt source the coefficient at an extracted prime
would be strictly negative, which is impossible. The remaining
discrepancy inequalities show that the model is klt.
\end{proof}

\begin{proposition}[Semiampleness on the reduced boundary]
\label{proj:prop:boundary-restriction}
Assume Assumption~\ref{proj:mmp:lower-models}. Let \((V,D)\) be a
projective \(\Q\)-factorial dlt \(n\)-fold with \(D\) reduced and
\(M=K_V+D\) nef. Then \(M|_D\) is semiample as a rational line
bundle on the reduced scheme \(D\).
\end{proposition}

\begin{proof}
The ambient variety \(V\) is klt. The reduced floor of a
\(\Q\)-factorial dlt pair is \(S_2\), and it has ordinary double
crossings in codimension one; its irreducible components are normal.
For the \(S_2\) assertion one can work locally with the cyclic
class cover of the \(\Q\)-Cartier divisor \(D\). This cover is
quasi-\'etale and klt, hence Cohen--Macaulay. The eigensheaf
\(\OO_V(-D)\), being a direct summand of its finite direct image,
is Cohen--Macaulay. The divisor sequence then shows that \(D\)
is Cohen--Macaulay.

Divisorial adjunction on the normalization
\(\coprod D_i\to D\) gives lc pairs
\((D_i,\operatorname{Diff}_{D_i}(D-D_i))\).
The different includes the conductor with coefficient one.
Removing that conductor contribution defines an effective boundary
on \(D\); together with \(D\) it is an slc pair whose adjoint
line bundle is \(M|_D\). The divisible residue identifications
hold in codimension one, with even powers eliminating residue
signs at double crossings, and extend by \(S_2\).

Each normalized adjoint is nef and semiample by
Assumption~\ref{proj:mmp:lower-models} and
Lemma~\ref{proj:mmp:comparison}.
Semiampleness descends from the normalization by the slc gluing
theorem \cite[Theorem~1.5, equivalently Theorem~4.3]
{FujinoGongyoPluricanonical}. This gives the required statement on the
whole reduced scheme, not merely on its individual components.
\end{proof}

\subsection{The special termination needed below}

\begin{proposition}[Special termination over a point]
\label{proj:prop:special-termination}
Assume Assumption~\ref{proj:mmp:lower-models}.
Let \((V,\Delta)\) be a projective \(\Q\)-factorial dlt
\(n\)-fold with rational boundary and pseudo-effective adjoint.
Choose an effective ample rational divisor \(A\) such that
\((V,\Delta+A)\) is dlt and \(K_V+\Delta+A\) is nef.
The LMMP for \(K_V+\Delta\) with scaling of \(A\) has special
termination: if the program is infinite, its flipping loci
are eventually disjoint from the round-down of the boundary.

In particular, if at every stage the adjoint is rationally linearly
equivalent to a signed divisor supported on that round-down, the
program terminates at a log minimal model.
\end{proposition}

\begin{proof}
In an infinite ample-scaling program the scaling limit is zero:
a positive limit is covered by the termination theorem for a
dlt pair with an ample summand in the scaling divisor
\cite[Theorem~1.9(ii), equivalently Theorem~4.1(ii)]
{Proj-Birkar2012Special}. Discard the finitely many divisorial
contractions and write the flips as \(V_i\dashrightarrow V_{i+1}/Z_i\),
with scaling numbers \(\lambda_i>0\) tending to zero.
Fix a component \(S_1\) of \(\lfloor\Delta_1\rfloor\), write
\(S_i\) for its normal strict transform, and let \(T_i\) be the
normalization of its image in \(Z_i\). Since the ambient contraction
is small, \(S_i\to T_i\) is projective birational. The standard
discrepancy argument makes \(S_i\dashrightarrow S_{i+1}\)
an isomorphism in codimension one after finitely many steps
\cite[proof of Lemma~3.6]{Proj-Birkar2010}.

Here are the precise auxiliary relative programs in that argument.
Take a small projective \(\Q\)-factorialization
\(p_i:S'_i\to S_i\), and put
\[
 D_i=K_{S'_i}+B_i=p_i^*((K_{V_i}+\Delta_i)|_{S_i}),
 \qquad C_i=p_i^*(A_i|_{S_i}).
\]
Adjunction gives a dlt pair \((S'_i,B_i)\), with \(B_i,C_i\geq0\),
and \((S'_i,B_i+\lambda_i C_i)\) is lc. The scaling divisor
has no floor component, so its restriction is effective.
The number \(\lambda_i\) is rational, being the ratio of
intersections of rational divisors on the contracted rational curve.
The globally nef divisor
\(L_i=K_{V_i}+\Delta_i+\lambda_i A_i\) descends rationally
linearly across \(V_i\to Z_i\). Indeed, for sufficiently small
rational \(\delta>0\), the pair
\((V_i,(1-\delta)\Delta_i)\) is klt and its negative adjoint
is ample over \(Z_i\). Relative basepoint freeness applies to a
Cartier multiple of \(L_i\); its numerical triviality and the
connected fibers give descent. Restricting and pulling back yields
\[
 D_i+\lambda_i C_i\sim_{\Q}0/T_i.
\]

By \cite[Theorem~1.1(3)]{Proj-Birkar2012Special}, a \(D_i\)-LMMP
over \(T_i\) with scaling of a fresh ample divisor terminates.
The theorem applies to the effective rational lc boundary
\(B_i+\lambda_i C_i\), with \(\lambda_i C_i\) rational Cartier
and the driving pair \(\Q\)-factorial dlt. Since the base map
is birational, the endpoint is a log minimal model, not a Mori
fiber space. Comparison with the relatively ample model
\(S_{i+1}\) identifies this endpoint with a small
\(\Q\)-factorialization \(S'_{i+1}\), as in
\cite[Remarks~2.9--2.10]{Proj-Birkar2012Special}.

These finite relative programs are genuine pieces of an absolute
program with scaling of the transforms of \(C_1\).
Throughout the \(i\)-th piece, \(D_i+\lambda_i C_i\) remains
the pullback of a nef divisor on \(T_i\), hence is globally nef.
Every contracted negative ray has positive \(C_i\)-degree,
so its global scaling threshold is exactly \(\lambda_i\).
The relative cone is a face of the absolute cone, cut out by
the pullback of an ample divisor on the projective base \(T_i\);
its extremal rays are therefore absolute extremal rays.
Thus the pieces concatenate as asserted in the cited remarks.
Rescaling the initial scaling divisor by \(\lambda_1\), if
necessary, makes its sum with the initial boundary lc and nef.

If the concatenation were infinite, its positive scaling numbers
would tend to zero. Its initial adjoint is pseudo-effective:
for each late scaling value, pull back the corresponding nef
scaled adjoint and use the effective comparison divisor for
the preceding nonpositive steps, then take the limit.
Assumption~\ref{proj:mmp:lower-models} therefore gives an absolute
log minimal model for that initial pair. The concatenation
terminates by \cite[Theorem~1.9(iii)]{Proj-Birkar2012Special}, since
its zero limit is never attained. The floor-component argument
of \cite[Lemma~3.6]{Proj-Birkar2010} now gives special termination.
No effective representative in dimension \(n\), nor any
arbitrary relative good-model assertion, has been used.

For the final assertion, a curve disjoint from the round-down
has degree zero against any signed divisor supported there.
It therefore cannot generate an adjoint-negative flipping ray.
Special termination rules out all sufficiently late flips.
There are only finitely many divisorial contractions, since
each lowers the Picard number, and pseudo-effectivity excludes
a Mori fiber space as the endpoint. Thus the endpoint is nef.
\end{proof}

\subsection{A uniform trivial-perturbation argument}

\begin{lemma}\label{proj:mmp:trivial-perturbation}
Let \((Z,\Delta)\) be a projective \(\Q\)-factorial dlt rational
pair, and suppose \(L=K_Z+\Delta\) is nef.
Set \(P=\Delta\) if the pair is klt and
\(P=\lfloor\Delta\rfloor\) otherwise.
Fix \(m>0\) such that \(mL\) is Cartier.
For every rational
\[
 0<\epsilon<\frac{1}{4nm+2},
\]
every step of an LMMP for \(L-\epsilon P\) is \(L\)-trivial.
On all its models, the transform of \(L\) is nef and its
multiple by the same integer \(m\) is Cartier.
\end{lemma}

\begin{proof}
Both \(\Delta-\epsilon P\) and \(\Delta-\tfrac12P\) are
effective klt boundaries. If \(R\) is a ray negative for
\(L-\epsilon P\), nefness of \(L\) shows that \(R\) is also
negative for \(L-\tfrac12P\).
Choose a rational curve \(C\) spanning this ray and satisfying
the length bound
\[
 -\bigl(L-\tfrac12P\bigr)\cdot C\leq 2n.
\]
Writing \(a=L\cdot C\geq 0\) and \(p=P\cdot C\), we have
\[
 p\leq 2a+4n,\qquad a<\epsilon p,
 \qquad
 (1-2\epsilon)a<4n\epsilon.
\]
Our choice of \(\epsilon\) gives \(a<1/m\).
Since \(ma\) is a nonnegative integer, \(a=0\).

For the driving klt contraction, the line bundle
\(\OO_Z(mL)\) is numerically trivial over the contraction base.
The line-bundle clause of the contraction theorem gives its
descent as a line bundle, with no change of \(m\)
\cite[Theorem~3.74(3)]{Proj-Fujino2017}.
One may also see the unchanged index directly from relative
basepoint freeness: two consecutive sufficiently large powers
descend, so their quotient descends.
For a flip, pull this descended line bundle back on the other
side. The transformed \(L\) is nef, and \(mL\) remains Cartier.

The driving boundary remains klt throughout its LMMP.
The transformed \(\Delta-\tfrac12P\) remains effective and
is smaller than the driving boundary, so it too is klt.
The same estimate and the same integer \(m\) apply inductively
at every subsequent step.
\end{proof}

\subsection{Completion of the good-model step}

\begin{proposition}[Inductive reduction to smooth nonvanishing]
\label{proj:prop:inductive-reduction}
Assume Assumption~\ref{proj:mmp:lower-models}. Assume in addition that
every smooth projective \(n\)-fold with pseudo-effective
canonical divisor has nonnegative Kodaira dimension.
Then every projective lc \(n\)-fold with real boundary and
pseudo-effective adjoint has a good log minimal model.
In particular, every nef rational lc adjoint in dimension \(n\)
is semiample.
\end{proposition}

\begin{proof}
Hashizume's reduction gives nonvanishing and log minimal models
for projective lc pairs with real boundary in dimensions at most
\(n\) \cite[Theorem~1.4]{Proj-Hashizume2018}.
Thus we may pass to a \(\Q\)-factorial dlt log minimal model.
Fix the support of its boundary, impose the rational affine
constraints that its coefficient-one components remain equal
to one, and take a sufficiently small rational polytope around
the given boundary within those constraints. On a fixed log
resolution witnessing dlt, the strict discrepancy inequalities
remain strict in this neighborhood; thus its boundaries remain
dlt. The rational-polytope theorem for nef adjoints
\cite[Remark~3.1 and Proposition~3.2(3)]{Proj-Birkar2011II}, applied
to all extremal rays and intersected with this neighborhood,
expresses the given real boundary in a finite rational simplex
of dlt boundaries with nef adjoints.
It is enough to prove semiampleness for these
rational adjoints. Their positive real combinations are
semiample, using the product of the associated contractions.
For a rational adjoint, real semiampleness can also be
rationalized: the finite linear equations expressing its
divisor and principal-divisor coefficients have rational data,
so a real solution with positive coefficients has a rational
solution with the same positivity.

We therefore work with a rational dlt pair \((Z,\Delta)\) and
\(L=K_Z+\Delta\) nef.
Real nonvanishing gives rational nonvanishing in this case:
retain the finitely many divisors and principal divisors in
an effective real representative and solve the resulting
rational linear system with nonnegative rational coefficients.
If \(\kappa(Z,L)\geq 1\), lower-dimensional good models give a
good minimal model by \cite[Lemma~3.5]{Proj-FujinoGongyoRings2014}.
Lemma~\ref{proj:mmp:comparison} transfers semiampleness back to \(Z\).
It remains to treat
\[
 \kappa(Z,L)=0.
\]

\paragraph{The non-pseudo-effective perturbation.}
Put \(P=\Delta\) in the klt case and
\(P=\lfloor\Delta\rfloor\) in the other case.
Suppose \(L-\epsilon P\) is not pseudo-effective for every
sufficiently small \(\epsilon>0\).
Choose rational \(\epsilon\) as in
Lemma~\ref{proj:mmp:trivial-perturbation}, and run the klt LMMP
with scaling to a Mori fiber space
\[
 (Z,\Delta-\epsilon P)\dashrightarrow
 (Z',\Delta'-\epsilon P')\xrightarrow{f} T.
\]
Existence and termination in the non-pseudo-effective case are
the established klt results of \cite{BCHM}.
The program is crepant for \(L\).
The line-bundle descent in
Lemma~\ref{proj:mmp:trivial-perturbation}, applied also to the final
contraction, gives a nef rational divisor \(A\) on \(T\) with
\[
 K_{Z'}+\Delta'\sim_{\Q}f^*A.
\]
The base has dimension less than \(n\).

If \((Z,\Delta)\) is klt, the transformed original pair
\((Z',\Delta')\) is klt as well. Ambro's descent theorem
\cite[Theorem~0.2]{Proj-Ambro2005} gives an effective rational
boundary \(\Delta_T\) such that \((T,\Delta_T)\) is klt and
\[
 A\sim_{\Q}K_T+\Delta_T.
\]
All its hypotheses hold: the total pair is globally klt and
effective, \(f\) is a projective contraction, and its adjoint
is rationally linearly pulled back. In particular, this use
requires no conjecture about semiampleness of a moduli divisor.
Assumption~\ref{proj:mmp:lower-models} and
Lemma~\ref{proj:mmp:comparison} make \(A\) semiample.

Otherwise \(P'\) is effective and relatively ample, since
\(K_{Z'}+\Delta'-\epsilon P'\) is relatively antiample.
Some component of the floor therefore dominates \(T\).
On a dlt blowup of the original lc pair \((Z',\Delta')\),
choose the strict transform \(S\) of such a component.
Adjunction produces an lc pair on \(S\) whose nef adjoint
is the pullback of \(A\).
It is semiample by the lower-dimensional hypothesis.
Semiampleness of a rational line bundle descends along a proper
surjection: use the equality of sections for its connected-fiber
Stein factor, and norms for the remaining finite morphism.
For the latter assertion, a basepoint-free system has a section
nonzero at every point of a chosen finite fiber; its norm is
nonzero at the point downstairs. Hence \(A\) is semiample
also in this case. Crepancy transfers the conclusion back to \(L\).

\paragraph{The klt pseudo-effective perturbation.}
We next settle the klt case when \(L-\epsilon\Delta\) is
pseudo-effective for some small rational \(\epsilon>0\).
Nonvanishing gives
\[
 L\sim_{\Q}G_0:=H_0+\epsilon\Delta,\qquad H_0\geq 0.
\]
Take a resolution \(p\colon W\to Z\) with reduced SNC divisor
\(D\) consisting of all exceptionals and the strict support of
\(H_0+\Delta\). We have
\[
 K_W+D=p^*L+R,\qquad R\geq 0.
\]
Here every component of \(D\) has positive coefficient in
\[
 G_W:=p^*G_0+R\sim_{\Q}K_W+D.
\]
Indeed, at strict boundary components the difference between
coefficient one and a klt boundary coefficient is positive,
and at exceptional components the coefficient contributed by
the adjoint comparison is the positive log discrepancy.
The pushforward \(p_*G_W\) is bounded coefficientwise by a
fixed multiple of \(G_0\). Section spaces inject under
birational pushforward, so
\[
 0\leq\kappa(W,K_W+D)\leq\kappa(Z,G_0)=0.
\]

Take a \(\Q\)-factorial dlt log minimal model
\((V,D_V)\) of \((W,D)\).
Its boundary is reduced. On a common resolution, the comparison
divisor in Lemma~\ref{proj:mmp:comparison} is exceptional over \(V\).
Pushing the pullback of \(G_W\) to \(V\) therefore gives
\[
 K_V+D_V\sim_{\Q}G_V=\sum_i b_iD_i,\qquad b_i>0,
 \qquad \Supp G_V=D_V.
\]
For completeness, positivity at a component extracted on \(V\)
also holds: its log discrepancy over \((W,D)\) is zero, its
center lies in \(D\), and the pullback of the full-support
effective \(G_W\) has positive order there.
There is no comparison-divisor coefficient at a prime retained
on \(V\).
Effectivity and exceptionality in the comparison preserve the
adjoint section ring, so the nef adjoint on \(V\) still has
Iitaka dimension zero.

The small-perturbation hypothesis of
Theorem~\ref{proj:thm:signed} is now explicit.
If \(D_V\neq 0\), choose rational \(c\) with
\(0<c<\min_i b_i\). Then
\[
 K_V+(1-c)D_V
 \sim_{\Q}\sum_i(b_i-c)D_i\geq 0.
\]
If \(D_V=0\), pseudo-effectivity is immediate.
Proposition~\ref{proj:prop:boundary-restriction} supplies the
required semiampleness on \(D_V\).
Theorem~\ref{proj:thm:signed} thus makes \(K_V+D_V\) abundant.
Its numerical dimension is zero, and its effective
full-support representative forces \(G_V=0\), hence \(D_V=0\).

On a common resolution of \(W\dashrightarrow V\), the effective
pullback of \(p^*G_0\) is now exceptional over \(V\).
It is nef because it is rationally linearly equivalent to the
pullback of \(L\). The negativity lemma forces this divisor
to vanish. Therefore \(G_0=0\) and \(L\sim_{\Q}0\).
Together with the preceding cases, this proves the klt
good-model assertion in dimension \(n\): for a non-nef
pseudo-effective klt adjoint, first take its klt log minimal
model and apply what was just proved.

\paragraph{The remaining non-klt case.}
Finally suppose the original pair is not klt and a small
floor perturbation is pseudo-effective. Choose rational
\(\epsilon>0\) so that both
\[
 L-\epsilon P,\qquad L-2\epsilon P,
 \qquad P=\lfloor\Delta\rfloor,
\]
are pseudo-effective klt adjoints.
Their Iitaka dimensions are zero: nonvanishing gives the lower
bound, and adding the effective perturbation bounds each by
\(\kappa(Z,L)=0\).
The klt result just proved gives good models for both.
Consequently their Nakayama numerical dimensions are zero.
Since
\[
 2(L-\epsilon P)
   =L+(L-2\epsilon P)
\]
and the last summand has an effective rational representative,
monotonicity and homogeneity of \(\kappa_\sigma\) give
\[
 \kappa_\sigma(Z,L)
 \leq\kappa_\sigma\bigl(Z,2(L-\epsilon P)\bigr)=0.
\]
The nef divisor \(L\) is therefore numerically trivial.
Nonvanishing is already available in this finishing step:
write \(L\sim_{\Q}E\geq 0\).
For an ample divisor \(H\), the equality
\(E\cdot H^{n-1}=0\) forces \(E=0\).
Thus \(L\sim_{\Q}0\), without any additional nonvanishing or
abundance premise.

All rational nef dlt adjoints are now semiample.
The polytope reduction gives the real-boundary good-model
assertion, and Lemma~\ref{proj:mmp:comparison} gives the final
statement for every nef rational lc adjoint.
\end{proof}
